% TEMPLATE-MILC-v3.tex - plantilla maestra para todas las guias MILC v3
% Ver scripts/generadores/build-guia-milc.py para la lista completa de
% claves esperadas. El script valida que no queden placeholders sin
% reemplazar antes de escribir el archivo final.

\documentclass[11pt,letterpaper]{article}
\usepackage[margin=1.75cm]{geometry}
\usepackage[table]{xcolor}
\usepackage{fontspec}
\usepackage[spanish]{babel}
\usepackage{tikz}
\usepackage[most]{tcolorbox}
\usepackage{tabularx}
\usepackage{array}
\usepackage{fancyhdr}
\usepackage{titlesec}
\usepackage{ragged2e}
\usepackage{enumitem}
\usepackage{microtype}

\setmainfont{Helvetica Neue}
\setsansfont{Helvetica Neue}
\setlength{\parindent}{0pt}
\setlength{\parskip}{6pt}
\hyphenpenalty=200
\exhyphenpenalty=200
\pretolerance=400
\tolerance=2000
\setlength{\emergencystretch}{1em}
\renewcommand{\arraystretch}{1.25}
\setlist{leftmargin=5mm,itemsep=1mm,topsep=1mm}
\newcolumntype{Y}{>{\RaggedRight\arraybackslash}X}

\definecolor{milcMagenta}{HTML}{E6007E}
\definecolor{milcVino}{HTML}{5A0038}
\definecolor{milcTurquesa}{HTML}{0093A5}
\definecolor{milcVerde}{HTML}{58A923}
\definecolor{milcMostaza}{HTML}{E5B400}
\definecolor{milcOcre}{HTML}{B86600}
\definecolor{milcNegro}{HTML}{241816}
\definecolor{milcGris}{HTML}{F7F7F4}
\definecolor{milcLinea}{HTML}{E8E2DD}
\definecolor{milcGradePrimary}{HTML}{0066FF}
\definecolor{milcGradeDark}{HTML}{003D99}
\definecolor{milcGradeSoft}{HTML}{D6E8FF}
\definecolor{milcGradeAccent}{HTML}{CFE7FF}
\definecolor{milcGradeText}{HTML}{FFFFFF}
\color{milcNegro}

\pagestyle{fancy}
\fancyhf{}
\lhead{\small Tecnología e Informática · Grado 6}
\rhead{\small Guía 17 · MILC v3}
\cfoot{\small\thepage}
\renewcommand{\headrulewidth}{0pt}

\newtcolorbox{titlebox}[1]{
  enhanced,colback=#1,colframe=#1,arc=7mm,boxrule=0pt,
  left=7mm,right=7mm,top=3mm,bottom=3mm,
  before skip=8pt,after skip=8pt,
  fontupper=\sffamily\bfseries\Large\color{white}
}

\newtcolorbox{softbox}[3]{
  enhanced,colback=#2,colframe=#1,borderline west={4pt}{0pt}{#1},
  arc=2mm,boxrule=.4pt,left=4mm,right=4mm,top=3mm,bottom=3mm,
  before skip=6pt,after skip=8pt,title={#3},coltitle=white,
  colbacktitle=#1,fonttitle=\sffamily\bfseries,fontupper=\RaggedRight,
  attach boxed title to top left={xshift=3mm,yshift=-2mm},
  boxed title style={arc=2mm, boxrule=0pt}
}

\newtcolorbox{drawbox}[2]{
  enhanced,colback=white,colframe=#1,arc=4mm,boxrule=1pt,
  left=5mm,right=5mm,top=4mm,bottom=4mm,before skip=8pt,after skip=8pt,
  title={#2},coltitle=white,colbacktitle=#1,fonttitle=\sffamily\bfseries,
  fontupper=\RaggedRight,
  attach boxed title to top left={xshift=4mm,yshift=-2mm},
  boxed title style={arc=3mm, boxrule=0pt}
}

\newtcolorbox{infoband}[1]{
  enhanced,colback=#1!8,colframe=#1!50,arc=2mm,boxrule=.5pt,
  left=4mm,right=4mm,top=2mm,bottom=2mm,before skip=4pt,after skip=4pt,
  fontupper=\small\RaggedRight
}

% Puente narrativo entre secciones (contrato editorial Conéctate).
% Da continuidad: "Acabas de X. Ahora vas a Y porque Z."
% Visualmente: banda izquierda en color del grado, texto en itálica pequeña.
\newtcolorbox{puentebox}{
  enhanced,colback=milcGradeSoft,colframe=milcGradeDark,
  borderline west={3pt}{0pt}{milcGradeDark},
  arc=0mm,boxrule=0pt,left=5mm,right=4mm,top=2.5mm,bottom=2.5mm,
  before skip=4pt,after skip=8pt,
  fontupper=\itshape\small\color{milcNegro}\RaggedRight
}
\newcommand{\puente}[1]{%
  \begin{puentebox}%
    \textcolor{milcGradeDark}{\textbf{\textrightarrow}}\hspace{2mm}#1%
  \end{puentebox}%
}

\newcommand{\linead}[1]{\noindent\color{milcLinea}\rule{#1}{0.5pt}\color{milcNegro}}
\newcommand{\respuesta}[1]{\par\vspace{2mm}\foreach \n in {1,...,#1}{\noindent\color{milcLinea}\rule{\linewidth}{0.5pt}\par\vspace{5mm}}\color{milcNegro}}
\newcommand{\checkbox}{\textcolor{milcGradeDark}{\fbox{\phantom{x}}}\hspace{5pt}}

\newcommand{\phasecard}[4]{
\begin{tcolorbox}[enhanced,colback=#3,colframe=#2,arc=3mm,boxrule=.5pt,height=3.05cm,valign=center,halign=center,left=1.5mm,right=1.5mm,top=2mm,bottom=2mm]
{\sffamily\bfseries\fontsize{22}{24}\selectfont\textcolor{#2}{#1}}\par
{\sffamily\bfseries\small #4}
\end{tcolorbox}
}

\begin{document}
\thispagestyle{empty}
\begin{tikzpicture}[remember picture,overlay]
  \fill[white] (current page.south west) rectangle (current page.north east);
  \path[fill=milcGradePrimary,rounded corners=16mm]
    ([xshift=.55cm,yshift=-.55cm]current page.north west) rectangle
    ([xshift=-.55cm,yshift=3.65cm]current page.south east);

  \node[anchor=north west,text=milcGradeText] at ([xshift=2.0cm,yshift=-1.85cm]current page.north west)
    {\sffamily\bfseries\fontsize{14}{16}\selectfont GRADO};
  \node[anchor=north west,text=milcGradeText,opacity=.82] at ([xshift=2.0cm,yshift=-2.75cm]current page.north west)
    {\sffamily\bfseries\fontsize{9}{11}\selectfont SERIE GUÍAS · TECNOLOGÍA E INFORMÁTICA};

  \begin{scope}[shift={([xshift=-3.05cm,yshift=-2.55cm]current page.north east)}]
    \fill[white,opacity=.80] (-.62,-.52) rectangle (.62,.52);
    \draw[milcGradeText,line width=1pt] (-.62,-.52) rectangle (.62,.52);
    \fill[milcGradeDark] (-.42,-.38) rectangle (-.22,.06);
    \fill[milcGradeAccent] (-.10,-.38) rectangle (.10,.24);
    \fill[milcGradePrimary!60!black] (.22,-.38) rectangle (.42,.38);
  \end{scope}

  \node[anchor=west,text=milcGradeText] at ([xshift=1.9cm,yshift=-12.85cm]current page.north west)
    {\sffamily\bfseries\fontsize{124}{124}\selectfont 6};
  \draw[milcGradeText,line width=1.7pt]
    ([xshift=4.52cm,yshift=-11.68cm]current page.north west) circle (.13cm);

  \node[anchor=west,text=milcGradeText,text width=8.7cm,align=left] at ([xshift=10.35cm,yshift=-11.6cm]current page.north west)
    {
     {\sffamily\bfseries\fontsize{19}{22}\selectfont Archivos y\\carpetas ---\\ordenar\\mi biblioteca\\digital}\\[4mm]
     {\sffamily\bfseries\fontsize{23}{26}\selectfont Guía 17-6-TIC}\\[2mm]
     {\sffamily\fontsize{13}{16}\selectfont Período 2 · Hardware y software}\\[5mm]
     {\sffamily\bfseries\fontsize{11}{13}\selectfont Institución Educativa Sor María Juliana}\\[1mm]
     {\sffamily\fontsize{10}{12}\selectfont Autor: Dr. Álvaro Cárdenas Orozco}
    };

  \path[fill=milcGradeSoft,rounded corners=7mm]
    ([xshift=1.35cm,yshift=.85cm]current page.south west) rectangle
    ([xshift=-1.35cm,yshift=4.25cm]current page.south east);
  \draw[milcGradeDark,line width=.65pt,rounded corners=7mm]
    ([xshift=1.35cm,yshift=.85cm]current page.south west) rectangle
    ([xshift=-1.35cm,yshift=4.25cm]current page.south east);
  \node[anchor=center,text=milcNegro,text width=17.0cm,align=center] at ([yshift=2.55cm]current page.south)
    {\sffamily\fontsize{10}{12}\selectfont Metodología MILC · Apertura ancestral · Pensamiento computacional · Triángulo Dussel-Estoicismo-Floridi\\
    \textcolor{milcGradeDark}{\bfseries Producto:} Una \textbf{estructura personal de carpetas} dibujada en árbol (jerarquía) en una hoja del cuaderno, con \textbf{mínimo 4 niveles de profundidad}, \textbf{15 carpetas} bien nombradas y \textbf{10 archivos reales} (con su extensión) ubicados en la carpeta correcta. Más \textbf{una tabla de las 10 extensiones más comunes} (.docx, .jpg, .pdf, .mp4, .mp3, etc.) con qué significan. Firma: \emph{``Mis archivos están ordenados como en una biblioteca, no como en un caos.''}.};
\end{tikzpicture}
\mbox{}
\newpage

\section*{Apertura: Archivos y carpetas --- ordenar mi biblioteca digital}

\begin{softbox}{milcGradeDark}{milcGradeSoft}{Saber ancestral}
\textbf{Doña Elena trabajó 30 años como archivera del juzgado de Cartago.} En esa oficina llegaban cada día 50 expedientes, cartas, peticiones, sentencias. Si se mezclaban, una persona podía pasar meses sin que le respondieran un trámite. Doña Elena tenía un sistema: \textbf{archivadores de madera con etiquetas}, dentro de cada archivador \textbf{cajones} por tipo de caso, dentro de cada cajón \textbf{carpetas} por mes, dentro de cada carpeta \textbf{hojas numeradas} en orden cronológico. Si llegaba un abogado preguntando por \emph{``el expediente Martínez del 15 de marzo de 1995''}, doña Elena lo encontraba en menos de 2 minutos. \textbf{Su saber no era memoria, era organización}: cada cosa en su lugar, cada lugar con su nombre. Los abuelos en las casas hacían algo parecido con cajas de zapatos llenas de fotos: una por familia, una por viaje, una por niño. Hoy ese mismo sistema vive en tu computador: se llama \textbf{estructura de carpetas y archivos}. Aprenderlo te ahorra años de búsqueda perdida.
\end{softbox}

\begin{softbox}{milcTurquesa}{milcGris}{Saber contemporáneo}
Un \textbf{archivo} es una unidad de información digital: un documento de Word, una foto, una canción, un video, un PDF. Cada archivo tiene un \textbf{nombre} (la parte que tú escribes) y una \textbf{extensión} (las 3-4 letras después del punto). La extensión le dice al computador qué tipo de archivo es: \texttt{.docx} = Word, \texttt{.jpg} = foto, \texttt{.mp4} = video, \texttt{.pdf} = documento de lectura. Una \textbf{carpeta} es una caja virtual que guarda archivos (y también puede guardar otras carpetas adentro). Las carpetas se organizan en \emph{jerarquía}: como un árbol con tronco y ramas. Aprender a organizar bien es uno de los \textbf{4 hábitos más útiles} que adquieres en 6° grado: te ahorra tiempo, te quita estrés y te hace ver profesional cuando entregas trabajos.
\end{softbox}

\begin{softbox}{milcMostaza}{milcGris}{Pregunta-puente}
\textit{Tu tía te pide que le envíes \emph{``la foto de tu cumpleaños del año pasado''}. Si abres tu galería del celular y encuentras \textbf{4.000 fotos sin organizar}, ¿cuánto tardarías en encontrarla? ¿Y si las tuvieras en una carpeta llamada \emph{``2025-mi-cumpleaños''}?}
\end{softbox}

\begin{softbox}{milcVerde}{milcGris}{Saber y saber hacer}
Al terminar la clase: (1) podrás \textbf{identificar} archivos por su extensión; (2) sabrás \textbf{explicar} la jerarquía de carpetas (árbol); (3) podrás \textbf{aplicar} 5 reglas de nombrado a tus carpetas; (4) habrás \textbf{creado} una estructura personal de organización con archivos ubicados.
\end{softbox}

\small
\begin{tabularx}{\textwidth}{>{\bfseries}p{3.0cm}Y}
\rowcolor{milcGradePrimary!12}Autor & Dr. Álvaro Cárdenas Orozco\\
DBA & Identifica componentes y sistemas tecnológicos en su entorno (MEN, T\&I 6°)\\
Referentes & Saber ancestral del relojero · Sistemas como cadena de oficios · Diagnóstico paso a paso\\
Producto final & Una \textbf{estructura personal de carpetas} dibujada en árbol (jerarquía) en una hoja del cuaderno, con \textbf{mínimo 4 niveles de profundidad}, \textbf{15 carpetas} bien nombradas y \textbf{10 archivos reales} (con su extensión) ubicados en la carpeta correcta. Más \textbf{una tabla de las 10 extensiones más comunes} (.docx, .jpg, .pdf, .mp4, .mp3, etc.) con qué significan. Firma: \emph{``Mis archivos están ordenados como en una biblioteca, no como en un caos.''}.\\
\end{tabularx}
\normalsize

\puente{Hoy haces 4 cosas: identificas 10 archivos con sus extensiones, aprendes a armar carpetas en árbol, escribes las 5 reglas de nombrado, y dibujas tu árbol personal.}

\begin{titlebox}{milcGradeDark}
Ruta MILC: así vamos a aprender
\end{titlebox}
\begin{tabularx}{\textwidth}{>{\centering\arraybackslash}X>{\centering\arraybackslash}X>{\centering\arraybackslash}X>{\centering\arraybackslash}X}
\phasecard{1}{milcVerde}{milcGris}{Escuta\\[-1mm]\small Observo, escucho y pregunto.} &
\phasecard{2}{milcTurquesa}{milcGris}{Sistematización\\[-1mm]\small Pensamiento computacional.} &
\phasecard{3}{milcMagenta}{milcGris}{Praxis\\[-1mm]\small Construyo y aplico.} &
\phasecard{4}{milcVino}{milcGris}{Evaluación\\[-1mm]\small Reflexión liberadora.}
\end{tabularx}


\puente{Empezamos viendo 10 nombres de archivos para descifrar qué son por su extensión.}

\begin{titlebox}{milcVerde}
Fase 1 · Escuta
\end{titlebox}

\begin{softbox}{milcVerde}{milcGris}{Escena para escuchar}
\textbf{Actividad 1 · IDENTIFICA --- ¿Qué archivo es?} (10 min · individual). En el cuaderno, lee estos 10 nombres de archivos y \textbf{adivina qué tipo de archivo es} y \textbf{para qué sirve}. La \emph{extensión} (las 3-4 letras después del punto) es la pista. \texttt{(1) tarea-historia.docx} · \texttt{(2) foto-perfil.jpg} · \texttt{(3) cancion-favorita.mp3} · \texttt{(4) video-cumpleaños.mp4} · \texttt{(5) presentacion-clase.pptx} · \texttt{(6) factura-internet.pdf} · \texttt{(7) hoja-calculo.xlsx} · \texttt{(8) carpeta-fotos.zip} · \texttt{(9) instalador-juego.exe} · \texttt{(10) notas.txt}. Si no conoces alguna extensión, intenta deducir por el contexto del nombre.
\end{softbox}

\checkbox Copié los 10 nombres con extensión.\\
\checkbox Al lado de cada uno, escribí qué tipo es (Word, foto, audio, video, etc.).\\
\checkbox Apliqué intuición; no busqué en internet.

\begin{infoband}{milcVerde}
\textbf{Lo que va al cuaderno (Actividad 1 · IDENTIFICA).} Título: «Actividad 1 --- 10 archivos por su extensión». \textbf{Formato:} lista numerada del 1 al 10 con archivo + tipo. \textbf{Extensión:} 10 líneas. \textbf{Sabes que terminaste cuando} tienes los 10 archivos clasificados por tu cuenta.
\end{infoband}

\puente{Después del juego, aprendes la estructura de árbol y las 5 reglas para nombrar bien.}

\begin{titlebox}{milcTurquesa}
Fase 2 · Sistematización con pensamiento computacional
\end{titlebox}

\begin{softbox}{milcTurquesa}{milcGris}{Cuatro pilares aplicados al tema}
Lo que adivinaste son las \textbf{10 extensiones más comunes} que vas a ver en tu vida digital. Ahora vas a aprender qué significa cada una, cómo armar una estructura de carpetas, y las 5 reglas para nombrarlas bien.
\end{softbox}

\small
\begin{tabularx}{\textwidth}{>{\bfseries\RaggedRight\arraybackslash}p{3.6cm}Y}
\rowcolor{milcTurquesa!90}\color{white}Pilar & \color{white}\bfseries Aplicación al tema\\
1. Descomposición & \textbf{Las 10 extensiones más comunes.} \texttt{.docx} = documento de Word. \texttt{.pdf} = documento para leer (Adobe Reader). \texttt{.txt} = texto plano sin formato. \texttt{.jpg} y \texttt{.png} = fotos e imágenes (jpg ocupa menos, png tiene mejor calidad). \texttt{.gif} = imagen animada. \texttt{.mp3} = audio (canciones). \texttt{.mp4} y \texttt{.avi} = video. \texttt{.zip} y \texttt{.rar} = archivo comprimido (varias cosas en una caja). \texttt{.xlsx} = hoja de cálculo (Excel). \texttt{.pptx} = presentación (PowerPoint). \texttt{.exe} = programa que se instala (¡cuidado, solo de fuentes confiables!).\\
2. Reconocimiento de patrones & \textbf{La estructura de carpetas en árbol.} Imagina un árbol invertido. \textbf{La raíz} es la carpeta más alta: tu \emph{Documentos} o tu \emph{Mi PC}. De ahí salen \textbf{ramas} (subcarpetas) por temas: \emph{Colegio}, \emph{Familia}, \emph{Juegos}, \emph{Fotos}. Dentro de cada rama hay otras ramas más pequeñas: dentro de \emph{Colegio} podría haber \emph{Tecnología}, \emph{Matemáticas}, \emph{Sociales}. Y dentro de \emph{Tecnología} podría haber \emph{Periodo-1}, \emph{Periodo-2}. \textbf{La profundidad ideal} para 6° grado es de \emph{3 a 5 niveles}: ni muy plano (todo en una carpeta) ni muy profundo (te pierdes haciendo clic).\\
3. Abstracción & \textbf{Las 5 reglas de nombrado para no enredarte.} (1) \textbf{Usa palabras claras}: \emph{``tarea-historia''} es mejor que \emph{``hjfdksjf''}. (2) \textbf{Sin espacios, usa guion o guion bajo}: \emph{``mi-cumpleaños''} es mejor que \emph{``mi cumpleaños''} (algunos programas se enredan con espacios). (3) \textbf{Sin acentos ni ñ en el nombre}: aunque hoy funcionan, en algunos sistemas se rompen. \emph{``cumpleanos''} es más seguro que \emph{``cumpleaños''}. (4) \textbf{Pon fecha cuando aplique}: \emph{``2025-cumple''} es mejor que \emph{``cumple''} (después tendrás varios cumples). (5) \textbf{Nombra para tu yo de un año después}: pregúntate \emph{``¿yo el año entrante sabré qué hay en esta carpeta solo leyendo el nombre?''}. Si la respuesta es no, mejora el nombre.\\
4. Algoritmo conceptual & \textbf{Errores frecuentes al organizar.} (1) \emph{``Guardo todo en el Escritorio''} --- termina como tu cuarto desordenado, no encuentras nada. (2) \emph{``Pongo nombres como tarea1, tarea2, tarea3''} --- después no sabes cuál era cuál. (3) \emph{``Borro la papelera cada día''} --- si te equivocas, no recuperas. (4) \emph{``Pongo todo en una sola carpeta''} --- pierdes 15 minutos para encontrar algo. (5) \emph{``No hago copia de seguridad''} --- el día que falle el disco, perdiste todo el año. Hacer copia en USB o en Drive cada semana te salva.\\
\end{tabularx}
\normalsize

\begin{drawbox}{milcTurquesa}{Las 5 reglas + las 10 extensiones}
\textbf{5 reglas de nombrado:} \\[1mm]
\textbf{1.} Palabras claras (no códigos raros). \\[1mm]
\textbf{2.} Sin espacios (usa guion). \\[1mm]
\textbf{3.} Sin acentos ni ñ. \\[1mm]
\textbf{4.} Pon fecha si aplica. \\[1mm]
\textbf{5.} Piensa en tu yo de un año después. \\[1mm]
\textbf{10 extensiones:} \\[1mm]
.docx · .pdf · .txt · .jpg/.png · .mp3 · .mp4 · .zip · .xlsx · .pptx · .exe
\end{drawbox}

\begin{softbox}{milcMostaza}{milcGris}{Errores comunes y su corrección}
\textbf{Mitos sobre archivos y carpetas que escucharás.} (1) \emph{``Si guardo todo en el escritorio, es más rápido''} --- Falso. Se vuelve un caos en 3 meses. (2) \emph{``El .exe siempre es seguro''} --- Falso. Los .exe pueden ser virus. Solo instala los de fuentes oficiales. (3) \emph{``Comprimir un archivo (.zip) lo daña''} --- Falso. Lo hace más pequeño sin perder información. (4) \emph{``Si borro algo del computador, también se borra del Drive''} --- Depende. Si está sincronizado, sí. Si solo es copia, no. Verifica antes. (5) \emph{``Los nombres largos son siempre malos''} --- Falso. Mejor un nombre largo claro que uno corto incomprensible.
\end{softbox}

\begin{infoband}{milcTurquesa}
\textbf{Lo que va al cuaderno (Actividad 2 · EXPLICA).} Título: «Actividad 2 --- Reglas de nombrado y tabla de extensiones». \textbf{Formato:} 2 bloques. Bloque A: lista de las 5 reglas con un ejemplo bueno y uno malo. Bloque B: tabla de 10 filas con \emph{extensión + qué es + ejemplo de archivo}. \textbf{Extensión:} 5 líneas + tabla de 10 filas. \textbf{Sabes que terminaste cuando} podrías nombrar correctamente cualquier carpeta que crees a partir de hoy.
\end{infoband}

\puente{Con todo claro, diseñas tu árbol personal de carpetas con archivos reales.}

\begin{titlebox}{milcMagenta}
Fase 3 · Praxis con pensamiento computacional
\end{titlebox}

\begin{softbox}{milcMagenta}{milcGris}{Construyendo el producto con los cuatro pilares}
\textbf{Actividad 3 · CREA --- Tu árbol personal de carpetas} (28 min · individual). En una hoja completa del cuaderno vas a dibujar \textbf{tu estructura de carpetas} (en forma de árbol) con mínimo 4 niveles, 15 carpetas y 10 archivos reales tuyos. Esto será tu mapa para organizar tu computador o el celular.
\end{softbox}

\small
\begin{tabularx}{\textwidth}{>{\bfseries\RaggedRight\arraybackslash}p{3.6cm}Y}
\rowcolor{milcMagenta!90}\color{white}Pilar & \color{white}\bfseries Aplicación al producto\\
Descomposición & \textbf{Paso 1 --- La raíz.} En la parte superior de la hoja escribe el título: \textbf{``Mi árbol de carpetas''} y debajo: \emph{``[Tu nombre], año 2026''}. En el centro arriba, dibuja un \textbf{ícono de carpeta grande} y escribe debajo: \texttt{Mis-documentos}. Esa es la raíz.\\
Patrones & \textbf{Paso 2 --- Nivel 2: las 4 áreas principales.} Desde la raíz, traza \textbf{4 líneas hacia abajo} que terminen en 4 carpetas más pequeñas: \texttt{Colegio}, \texttt{Familia}, \texttt{Personal}, \texttt{Descargas}. Estas son las 4 áreas grandes de tu vida digital. Si tienes otras (por ejemplo, \emph{Música} o \emph{Juegos}), agrégalas.\\
Abstracción & \textbf{Paso 3 --- Nivel 3: subcarpetas dentro de cada área.} De \texttt{Colegio} traza líneas a: \texttt{Tecnologia}, \texttt{Matematicas}, \texttt{Sociales}. De \texttt{Familia} traza líneas a: \texttt{Fotos}, \texttt{Videos}. De \texttt{Personal} traza líneas a: \texttt{Cuaderno-digital}, \texttt{Apuntes}. De \texttt{Descargas} traza líneas a: \texttt{Pendientes-revisar}, \texttt{Archivados}.\\
Algoritmo de armado & \textbf{Paso 4 --- Nivel 4: subcarpetas más específicas.} De \texttt{Tecnologia} (Colegio) traza a: \texttt{Periodo-1}, \texttt{Periodo-2}, \texttt{Periodo-3}. De \texttt{Fotos} (Familia) traza a: \texttt{2024}, \texttt{2025}, \texttt{2026}. De \texttt{Pendientes-revisar} (Descargas) traza a: \texttt{Esta-semana}. Aquí ya tienes \textbf{4 niveles de profundidad}.\\
\end{tabularx}
\normalsize

\begin{drawbox}{milcMagenta}{Antes de mostrar tu árbol}
\checkbox La raíz (Mis-documentos) está arriba con su ícono. \\[1mm]
\checkbox Hay 4 áreas en el nivel 2 (Colegio, Familia, Personal, Descargas o similares). \\[1mm]
\checkbox Cada área tiene subcarpetas en nivel 3. \\[1mm]
\checkbox Hay al menos 4 niveles de profundidad. \\[1mm]
\checkbox Hay 10 archivos reales tuyos colocados en su carpeta correcta. \\[1mm]
\checkbox Está firmado y la frase de cierre está subrayada.
\end{drawbox}

\begin{softbox}{milcMagenta}{milcGris}{Plantilla de trabajo}
\textbf{Estructura del árbol.} \textit{Nivel 1 (raíz):} Mis-documentos. \textit{Nivel 2:} Colegio, Familia, Personal, Descargas. \textit{Nivel 3:} Tecnologia, Matematicas, Fotos, Videos, etc. \textit{Nivel 4:} Periodo-1, 2024, etc. \textit{Hojas:} archivos con extensión. \textit{Esquina:} firma + frase subrayada.
\end{softbox}

\begin{infoband}{milcMagenta}
\textbf{Lo que va al cuaderno (Actividad 3 · CREA).} Título: «Actividad 3 --- Mi árbol personal de carpetas». \textbf{Formato:} árbol jerárquico con 4 niveles, 15 carpetas, 10 archivos reales. \textbf{Extensión:} 1 hoja entera. \textbf{Sabes que terminaste cuando} cualquier archivo nuevo que crees el próximo mes \emph{sabrá dónde vivir} sin que tengas que pensarlo.
\end{infoband}

\puente{El producto es ese árbol firmado + tabla de extensiones.}

\begin{titlebox}{milcMostaza}
Producto de la sesión
\end{titlebox}
\textbf{Árbol personal de carpetas (4 niveles, 15 carpetas, 10 archivos) · firmado}.
\begin{softbox}{milcMostaza}{milcGris}{Criterios de calidad}
(1) El árbol tiene una sola raíz en el nivel superior. (2) Hay 4 áreas en el nivel 2. (3) Hay subcarpetas en nivel 3 y 4. (4) 10 archivos están ubicados en la carpeta correcta con su extensión. (5) Está firmado y la frase de cierre está visible.
\end{softbox}

\puente{Antes de salir, verificas que tu árbol tiene 4 niveles y que cada archivo está en la carpeta correcta.}

\begin{titlebox}{milcVino}
Fase 4 · Evaluación liberadora — cinco dimensiones
\end{titlebox}

\begin{softbox}{milcVino}{milcGris}{Sentido de la evaluación}
Evaluar no es solo revisar una nota. Es reconocer qué aprendí, qué cambió en mí, qué emociones manejé, cómo me relaciono con la ciudad y la comunidad, y qué le dejaré a quien viene detrás.
\end{softbox}

\textbf{Mi reflexión liberadora en cinco dimensiones.} Completa cada frase iniciada:

\small
\begin{tabularx}{\textwidth}{>{\bfseries\RaggedRight\arraybackslash}p{3.4cm}Y>{\RaggedRight\arraybackslash}p{4.6cm}}
\rowcolor{milcVino}\color{white}Dimensión & \color{white}\bfseries Mi respuesta (frame starter) & \color{white}\bfseries Algo que descubrí\\
1. Personal & Lo que más cambió en mí fue \linead{6.5cm} & \linead{4.2cm}\\
2. Emocional & La emoción más fuerte fue \linead{2.5cm} y la regulé \linead{3cm} & \linead{4.2cm}\\
3. Ciudadana & En democracia esto importa porque \linead{6cm} & \linead{4.2cm}\\
4. Local (Cartago) & En mi barrio sería útil para \linead{6.5cm} & \linead{4.2cm}\\
5. Intergeneracional & A mi abuela/abuelo le diría \linead{6.5cm} & \linead{4.2cm}\\
\end{tabularx}
\normalsize

\begin{infoband}{milcVino}
\textbf{Para la próxima sustentación.} Vuelve a esta tabla en seis meses. Verás que las respuestas cambian — no porque lo aprendido cambie, sino porque tú cambias. La evaluación liberadora es una conversación contigo a través del tiempo.
\end{infoband}


\puente{Cerramos con 3 voces que te ayudan a entender por qué ordenar es un acto que te cuida en el futuro.}

\begin{titlebox}{milcGradeDark}
Triángulo de pensamiento · cierre filosófico
\end{titlebox}

\begin{softbox}{milcVino}{milcGris}{Dussel · lente del nosotros}
\textit{````El orden no es una virtud burguesa: es un acto de cuidado hacia el propio trabajo.''''}

Algunas personas creen que \emph{``ser ordenado es cosa de obsesivos''}. Es al revés: ser ordenado es \textbf{respetar tu propio trabajo}. Cuando ordenas tus archivos, le estás diciendo a tu yo de mañana: \emph{``el trabajo que hago hoy importa, lo cuido para que mañana lo encuentre''}. Doña Elena, la archivera, no era obsesiva: era profesional. Tú con tus archivos haces lo mismo.

\textbf{¿Cuido mi trabajo digital como cuido las cosas que valoro físicamente?}
\end{softbox}
\linead{\linewidth}\\[1mm]
\linead{\linewidth}

\begin{softbox}{milcMostaza}{milcGris}{Estoicismo · Marco Aurelio (emperador romano que vivía con disciplina diaria) · cuidado interior}
\textit{````Cada cosa en su lugar; cada lugar con su propósito. Eso da paz al espíritu.''''}

Cuando todo está en su lugar, \textbf{no tienes que pensar dónde está}. Te concentras en hacer las cosas, no en buscarlas. \emph{El orden ahorra energía mental}. Cuando todo está revuelto, gastas atención en buscar; cuando todo está organizado, gastas atención en crear. La diferencia, sumada a lo largo del año, son días enteros de tu vida.

\textbf{¿Cuánto tiempo gasto buscando archivos perdidos cada semana? ¿Qué haría con ese tiempo si lo ahorrara?}
\end{softbox}
\linead{\linewidth}\\[1mm]
\linead{\linewidth}

\begin{softbox}{milcTurquesa}{milcGris}{Floridi · lente de la infoesfera}
\textit{````Quien organiza la información, gobierna su vida. Quien no, vive a merced del caos digital.''''}

Hoy producimos más información que nunca: \emph{fotos, videos, documentos, audios, descargas}. \textbf{Quien sabe organizarla está en ventaja sobre quien no}: encuentra rápido, comparte cuando le piden, no pierde lo importante. Aprender a organizar archivos no es lujo, es habilidad de supervivencia digital. Tu árbol de hoy es el primer paso para ser de los que organizan.

\textbf{¿Estoy gobernando mi información o me está gobernando ella a mí (en forma de caos)?}
\end{softbox}
\linead{\linewidth}\\[1mm]
\linead{\linewidth}

\begin{titlebox}{milcVino}
Compromiso final
\end{titlebox}

\small
\begin{tabularx}{\textwidth}{>{\RaggedRight\arraybackslash}p{3.0cm}YYY}
\rowcolor{milcVino}\color{white}\bfseries Criterio & \color{white}\bfseries Lo logré & \color{white}\bfseries En proceso & \color{white}\bfseries Necesito apoyo\\
Comprensión del tema & Explico ideas centrales con mis palabras. & Reconozco ideas pero falta precisión. & Necesito repasar conceptos base.\\
Pensamiento computacional & Apliqué los 4 pilares al diseñar mi producto. & Apliqué algunos pilares. & Necesito releer la fase 2.\\
Aplicación y producto & Mi producto cumple los criterios y se puede revisar. & Producto funcional con ajustes pendientes. & Aún en construcción.\\
Reflexión 5 dimensiones & Respondí cada dimensión con honestidad. & Respondí algunas con detalle. & Mis respuestas son muy generales.\\
Triángulo de pensamiento & Conecté las 3 voces con mi trabajo real. & Conecté al menos una. & No relaciono las voces con mi proceso.\\
\end{tabularx}
\normalsize

\textbf{Hoy entendí que:}\\[1mm]
\linead{\linewidth}\\[1mm]
\linead{\linewidth}

\textbf{Mi compromiso para la próxima sesión es:}\\[1mm]
\linead{\linewidth}\\[1mm]
\linead{\linewidth}

\begin{softbox}{milcMostaza}{milcGris}{Cita de despedida · Marco Aurelio}
\textit{``El obstáculo en el camino se vuelve el camino. Lo que estorba al camino se vuelve camino.''}\\[1mm]
Cada error de hoy, bien observado, se convierte en pieza del camino que pisarás mañana.
\end{softbox}

\small
\begin{tabularx}{\textwidth}{>{\bfseries}p{3.6cm}Y}
\rowcolor{milcGradeDark!12}Nombre completo: & \linead{10cm}\\
Fecha: & \linead{4cm} \quad Curso/grupo: \linead{4cm}\\
Mi firma: & \linead{9cm}\\
\end{tabularx}
\normalsize

\begin{infoband}{milcGradeDark}
\textbf{Para la próxima sesión.} Esta semana, en el computador o celular, crea 3 carpetas siguiendo las 5 reglas de nombrado. Sube 5 archivos a las carpetas correctas. La próxima clase aprendes \textbf{mantenimiento del computador}: cómo cuidar el equipo para que dure años, no meses. Hoy organizaste lo de adentro; la próxima ves cómo cuidar el equipo por dentro y por fuera.
\end{infoband}

\end{document}
